% GENERATED FILE. Edit canonical lessons, then run npm run generate:books.
% canonical-source-hash: fnv1a64:6ed78ffa23df6628
% canonical-lessons: GU-C535-sukan, GU-C535-vahem, GU-C535-sanket, GU-C535-grah, GU-C535-kalvu

\chapter{Signs and Omens}
\label{ch:gu-signs-and-omens}
\hlchaptermodality{535}

\begin{chapteropening}
\textbf{By the end of this chapter:} \emph{I can talk about omens, signs and superstitions.}

Complete the last lesson of chapter 535: I can talk about omens, signs and superstitions.
\end{chapteropening}

\section[\texorpdfstring{\gu{શુકન} (śukan)}{શુકન (śukan)}]{\gu{શુકન} (śukan) — an omen}

\label{lesson:GU-C535-sukan}



\begin{quote}
Before the new one: say the Gujarati for a lottery, then the Gujarati for a chance meeting.
\end{quote}

\subsection*{You'll want to know: \gu{શુકન}}
\textbf{\gu{શુકન}} — \emph{śukan} — \textquotedblleft{}an omen\textquotedblright{}.

\begin{quote}
\gu{જો} \gu{સવારે} \gu{ગાય} \gu{સામે} \gu{આવશે} \gu{તો} \gu{દાદી} \gu{તેને} \gu{સારું} \gu{શુકન} \gu{માનશે}.
\end{quote}

\begin{quote}
\emph{jo savāre gāy sāme āvśe to dādī tene sārũ śukan mānśe.}
\end{quote}

\begin{quote}
If a cow comes towards her in the morning, my grandmother will take it as a good omen.
\end{quote}

A neuter noun.

\begin{cousinweb}
\textbf{\gu{શુકન}} (\emph{śukan}) goes back to Sanskrit \emph{śakuna}, \textquotedblleft{}a bird\textquotedblright{}, especially a bird of good or bad omen (Monier-Williams).
\end{cousinweb}

\begin{grammarlens}[title={What you've built}]
One more word for saying what happens if.
\end{grammarlens}

\subsection*{Your turn}
\begin{itemize}
\raggedright
  \item \emph{Say it:} \emph{śukan}
  \item \emph{Say it:} \emph{śukan}, once more
  \item \emph{Say it:} \emph{jo savāre gāy sāme āvśe to dādī tene sārũ śukan mānśe.} — If a cow comes towards her in the morning, my grandmother will take it as a good omen.
  \item \emph{Recall:} say the Gujarati for a lottery, then the Gujarati for a chance meeting, then say \emph{śukan} again
\end{itemize}

\begin{tcolorbox}[breakable,colback=teal!4,colframe=teal!35!black,title={Before you move on}]
What does \gu{શુકન} mean? (\textquotedblleft{}An omen\textquotedblright{}.) Say it once more.
\end{tcolorbox}
\section[\texorpdfstring{\gu{વહેમ} (vahem)}{વહેમ (vahem)}]{\gu{વહેમ} (vahem) — a superstition}

\label{lesson:GU-C535-vahem}



\begin{quote}
Before the new one: say the Gujarati for a chance meeting, then the Gujarati for an omen.
\end{quote}

\subsection*{You'll want to know: \gu{વહેમ}}
\textbf{\gu{વહેમ}} — \emph{vahem} — \textquotedblleft{}a superstition\textquotedblright{}.

\begin{quote}
\gu{જો} \gu{તમે} \gu{આ} \gu{વહેમ} \gu{છોડી} \gu{દેશો} \gu{તો} \gu{તમારી} \gu{ચિંતા} \gu{ઓછી} \gu{થશે}.
\end{quote}

\begin{quote}
\emph{jo tame ā vahem choḍī deśo to tamārī cintā ochī thaśe.}
\end{quote}

\begin{quote}
If you give up this superstition, you'll worry less.
\end{quote}

A masculine noun.

\begin{cousinweb}
\textbf{\gu{વહેમ}} (\emph{vahem}) is usually traced to Arabic \emph{wahm}, \textquotedblleft{}a fancy, a groundless idea\textquotedblright{}, through Persian.
\end{cousinweb}

\begin{grammarlens}[title={What you've built}]
One more word for saying what happens if.
\end{grammarlens}

\subsection*{Your turn}
\begin{itemize}
\raggedright
  \item \emph{Say it:} \emph{vahem}
  \item \emph{Say it:} \emph{vahem}, once more
  \item \emph{Say it:} \emph{jo tame ā vahem choḍī deśo to tamārī cintā ochī thaśe.} — If you give up this superstition, you'll worry less.
  \item \emph{Recall:} say the Gujarati for a chance meeting, then the Gujarati for an omen, then say \emph{vahem} again
\end{itemize}

\begin{tcolorbox}[breakable,colback=teal!4,colframe=teal!35!black,title={Before you move on}]
What does \gu{વહેમ} mean? (\textquotedblleft{}A superstition\textquotedblright{}.) Say it once more.
\end{tcolorbox}
\section[\texorpdfstring{\gu{સંકેત} (sanket)}{સંકેત (sanket)}]{\gu{સંકેત} (sanket) — a hint}

\label{lesson:GU-C535-sanket}



\begin{quote}
Before the new one: say the Gujarati for an omen, then the Gujarati for a superstition.
\end{quote}

\subsection*{You'll want to know: \gu{સંકેત}}
\textbf{\gu{સંકેત}} — \emph{sanket} — \textquotedblleft{}a hint\textquotedblright{}.

\begin{quote}
\gu{જો} \gu{આકાશમાં} \gu{કાળાં} \gu{વાદળ} \gu{હશે} \gu{તો} \gu{તે} \gu{વરસાદનો} \gu{સંકેત} \gu{છે}.
\end{quote}

\begin{quote}
\emph{jo ākāśmā̃ kāḷā̃ vādaḷ haśe to te varsādno sanket che.}
\end{quote}

\begin{quote}
If there are black clouds in the sky, that is a sign of rain.
\end{quote}

A masculine noun.

\begin{cousinweb}
\textbf{\gu{સંકેત}} (\emph{sanket}) is a tatsama, taken from Sanskrit \emph{saṃketa}, \textquotedblleft{}an agreement; a sign, a signal\textquotedblright{} (Monier-Williams).
\end{cousinweb}

\begin{grammarlens}[title={What you've built}]
One more word for saying what happens if.
\end{grammarlens}

\subsection*{Your turn}
\begin{itemize}
\raggedright
  \item \emph{Say it:} \emph{sanket}
  \item \emph{Say it:} \emph{sanket}, once more
  \item \emph{Say it:} \emph{jo ākāśmā̃ kāḷā̃ vādaḷ haśe to te varsādno sanket che.} — If there are black clouds in the sky, that is a sign of rain.
  \item \emph{Recall:} say the Gujarati for an omen, then the Gujarati for a superstition, then say \emph{sanket} again
\end{itemize}

\begin{tcolorbox}[breakable,colback=teal!4,colframe=teal!35!black,title={Before you move on}]
What does \gu{સંકેત} mean? (\textquotedblleft{}A hint\textquotedblright{}.) Say it once more.
\end{tcolorbox}
\section[\texorpdfstring{\gu{ગ્રહ} (grah)}{ગ્રહ (grah)}]{\gu{ગ્રહ} (grah) — a planet (in astrology)}

\label{lesson:GU-C535-grah}



\begin{quote}
Before the new one: say the Gujarati for a superstition, then the Gujarati for a hint.
\end{quote}

\subsection*{You'll want to know: \gu{ગ્રહ}}
\textbf{\gu{ગ્રહ}} — \emph{grah} — \textquotedblleft{}a planet (in astrology)\textquotedblright{}.

\begin{quote}
\gu{જો} \gu{ગ્રહો} \gu{સારા} \gu{હશે} \gu{તો} \gu{દાદા} \gu{આ} \gu{મહિને} \gu{નવી} \gu{દુકાન} \gu{ખોલશે}.
\end{quote}

\begin{quote}
\emph{jo graho sārā haśe to dādā ā mahine navī dukān kholśe.}
\end{quote}

\begin{quote}
If the planets are favourable, my grandfather will open a new shop this month.
\end{quote}

A masculine noun.

\begin{cousinweb}
\textbf{\gu{ગ્રહ}} (\emph{grah}) is a tatsama, taken from Sanskrit \emph{graha}, \textquotedblleft{}a seizer\textquotedblright{}, and so \textquotedblleft{}a planet\textquotedblright{}, one held to seize on a person's fortunes (Monier-Williams).
\end{cousinweb}

\begin{grammarlens}[title={What you've built}]
One more word for saying what happens if.
\end{grammarlens}

\subsection*{Your turn}
\begin{itemize}
\raggedright
  \item \emph{Say it:} \emph{grah}
  \item \emph{Say it:} \emph{grah}, once more
  \item \emph{Say it:} \emph{jo graho sārā haśe to dādā ā mahine navī dukān kholśe.} — If the planets are favourable, my grandfather will open a new shop this month.
  \item \emph{Recall:} say the Gujarati for a superstition, then the Gujarati for a hint, then say \emph{grah} again
\end{itemize}

\begin{tcolorbox}[breakable,colback=teal!4,colframe=teal!35!black,title={Before you move on}]
What does \gu{ગ્રહ} mean? (\textquotedblleft{}A planet (in astrology)\textquotedblright{}.) Say it once more.
\end{tcolorbox}
\section[\texorpdfstring{\gu{કળવું} (kaḷvũ)}{કળવું (kaḷvũ)}]{\gu{કળવું} (kaḷvũ) — to make out (to guess)}

\label{lesson:GU-C535-kalvu}



\begin{quote}
Before the new one: say the Gujarati for a hint, then the Gujarati for a planet.
\end{quote}

\subsection*{You'll want to know: \gu{કળવું}}
\textbf{\gu{કળવું}} — \emph{kaḷvũ} — \textquotedblleft{}to make out (to guess)\textquotedblright{}.

\begin{quote}
\gu{જો} \gu{હું} \gu{તેની} \gu{વાત} \gu{પહેલાં} \gu{કળત} \gu{તો} \gu{હું} \gu{ત્યાં} \gu{ન} \gu{જાત}.
\end{quote}

\begin{quote}
\emph{jo hũ tenī vāt pahelā̃ kaḷat to hũ tyā̃ na jāt.}
\end{quote}

\begin{quote}
If I had made out what he meant earlier, I wouldn't have gone there.
\end{quote}

Past, with \gu{તેણે} (he or she …): \gu{તેણે} \gu{વાત} \gu{કળી}, \emph{teṇe vāt kaḷī}.

\begin{cousinweb}
\textbf{\gu{કળવું}} (\emph{kaḷvũ}) is inherited from Sanskrit \emph{kaláyati}, \textquotedblleft{}notices, considers\textquotedblright{}, and CDIAL lists the Gujarati verb, \textquotedblleft{}to understand, guess\textquotedblright{} (CDIAL 2918).
\end{cousinweb}

\begin{grammarlens}[title={What you've built}]
One more word for saying what happens if.
\end{grammarlens}

\subsection*{Your turn}
\begin{itemize}
\raggedright
  \item \emph{Say it:} \emph{kaḷvũ}
  \item \emph{Say it:} \emph{kaḷvũ}, once more
  \item \emph{Say it:} \emph{jo hũ tenī vāt pahelā̃ kaḷat to hũ tyā̃ na jāt.} — If I had made out what he meant earlier, I wouldn't have gone there.
  \item \emph{Recall:} say the Gujarati for a hint, then the Gujarati for a planet, then say \emph{kaḷvũ} again
\end{itemize}

\begin{tcolorbox}[breakable,colback=teal!4,colframe=teal!35!black,title={Before you move on}]
What does \gu{કળવું} mean? (\textquotedblleft{}To make out (to guess)\textquotedblright{}.) Say it once more.
\end{tcolorbox}
